%======================= ABSTRACT =======================
\newpage
\begin{center}\textbf{ABSTRACT}\end{center}
\addcontentsline{toc}{section}{ABSTRACT}
\vspace{0.4cm}

\noindent
In Brazil, home automation hardware is already cheap --- generic Wi-Fi bulbs and
plugs from the Tuya ecosystem cost between R\$~30 and R\$~60 ---, yet turning that
hardware into usable automation remains out of reach for the typical household due to
three chained barriers: the \textit{commissioning} cost (proprietary apps, cloud
accounts and English-only developer portals), structural cloud dependency, and privacy
risks. This work presents DOMUS, a proof of concept of a \textit{local-first} home
automation hub, controllable by Brazilian Portuguese voice processed on the hub
itself, whose thesis is that democratization depends less on hardware price and more
on lowering the commissioning barrier (``cloud once, local forever''). The system was
built with Node.js/NestJS and the \textit{Adapter} pattern to isolate protocol
fragmentation, with speech recognition via \textit{Whisper} running on the hub. As
results, local control of the TP-Link Tapo~P110 device was validated over the KLAP
protocol (latencies of 201--332~ms) and the voice \textit{pipeline} reached 87.7\%
intent accuracy; total cost, word error rate (WER) and usability (SUS) are reported in
Chapter~\ref{cap:resultados}. We conclude for the technical feasibility of local
Portuguese voice control without permanent cloud dependency; full validation of the
commissioning barrier for Tuya devices remains ongoing work.

\vspace{0.5cm}
\noindent\textbf{Keywords:} home automation; \textit{local-first}; commissioning;
speech recognition; privacy; Internet of Things.
\newpage
